
\documentclass[10pt,a4paper]{article}

\usepackage[utf8]{inputenc}
\usepackage[T1]{fontenc}
\usepackage[portuguese]{babel}
\usepackage{lmodern,amsmath,amssymb,graphicx,geometry,float,caption,subcaption,microtype,xcolor}

\geometry{margin=1.8cm,top=1.7cm,bottom=1.7cm}
\setlength{\parindent}{0pt}
\setlength{\parskip}{5pt}
\captionsetup{font=small}

\title{
    \vspace{-1.3cm}
    \textbf{Projeto 1 - Diagramas de Voronoi}\\
    \large Implementação interativa com WebGL2 e Fragment Shader
}

\author{
    \textbf{Elementos do grupo:}\\
    71003 Guilherme Martins | 65745 Rodrigo Santos\\
    \small Computação Gráfica e Interfaces
}

\date{Outubro de 2026}

\begin{document}

\maketitle
\vspace{-0.7cm}

\section{Introdução}

Um diagrama de Voronoi divide o plano em regiões, designadas \emph{células}, associadas a um conjunto de pontos de referência (\emph{sites}). Cada célula contém os pontos do plano cuja distância ao respetivo site é menor ou igual à distância aos restantes.

Para um conjunto $S=\{p_1,\ldots,p_n\}$, a célula de $p_i$ é

\begin{equation}
V(p_i)=\{x\in\mathbb{R}^2:\ d(x,p_i)\leq d(x,p_j),\ \forall j\neq i\}.
\end{equation}

Neste projeto foi desenvolvida uma aplicação interativa em WebGL2 que representa estas regiões através de pesquisa no \emph{fragment shader}. O utilizador pode adicionar, mover e remover sites, alterar a vista e selecionar três métricas de distância.

As figuras apresentam resultados obtidos com três sites e diferentes configurações de visualização.

\section{Arquitetura e cálculo no fragment shader}

A aplicação desenha um quadrilátero que ocupa toda a área do \emph{canvas}. Para cada fragmento, o shader converte a sua posição para coordenadas do mundo, tendo em conta o centro da câmara, a escala e a proporção da janela.

Em seguida, percorre os sites ativos, calcula a distância a cada um e identifica o de menor distância. A cor final provém da cor associada a esse site.

A implementação suporta até 32 sites por meio de \emph{uniform arrays} com posições e cores.

A intensidade da cor diminui suavemente com a distância ao site, utilizando um fator exponencial

\[
I=\exp(-\varepsilon d_1)
\]

em que $d_1$ é a menor distância calculada.

Os marcadores são representados por pequenos anéis, permitindo identificar visualmente os pontos de referência.

\section{Métricas de distância}

Para $p=(p_x,p_y)$ e $q=(q_x,q_y)$, foram implementadas as seguintes funções:

\begin{align}
d_E(p,q)&=\sqrt{(p_x-q_x)^2+(p_y-q_y)^2},\\
d_M(p,q)&=|p_x-q_x|+|p_y-q_y|,\\
d_C(p,q)&=\max\bigl(|p_x-q_x|,|p_y-q_y|\bigr).
\end{align}

A seleção é efetuada pelas teclas \texttt{1}, \texttt{2} e \texttt{3}, respetivamente.

As mesmas posições dos sites originam fronteiras diferentes, uma vez que o conceito de proximidade depende da métrica escolhida.

\begin{figure}[H]
    \centering
    \includegraphics[width=.59\linewidth]{Screenshot 2026-10-24 at 21.48.27.png}
    \caption{Diagrama com três sites, usando distância Euclidiana e bordas desativadas.}
    \label{fig:eucl}
\end{figure}

\section{Deteção e visualização das bordas}

Para evidenciar a separação entre células, o shader procura as duas menores distâncias a cada fragmento, $d_1$ e $d_2$.

Uma região é assinalada como fronteira quando

\begin{equation}
|d_1-d_2|<\varepsilon.
\end{equation}

Na implementação apresentada, foi usado $\varepsilon=0.1$ em unidades do mundo, e a tecla \texttt{b} ativa ou desativa as linhas brancas.

Este valor fixo é uma limitação: a espessura aparente das bordas varia com o zoom, pelo que uma melhoria futura seria adaptar o limiar à escala e à resolução do canvas.

\begin{figure}[H]
    \centering
    \includegraphics[width=.59\linewidth]{Screenshot 2026-10-24 at 21.48.30.png}
    \caption{Mesma configuração Euclidiana, agora com a deteção de bordas ativada.}
    \label{fig:borders}
\end{figure}

\section{Resultados e comparação visual}

As Figuras~\ref{fig:man} e~\ref{fig:cheb} mostram que a alteração da métrica modifica a forma das regiões sem exigir a movimentação dos sites.

Na distância Manhattan, os segmentos podem apresentar orientações horizontais, verticais e diagonais; na Chebyshev, observam-se também fronteiras com orientações características da norma do máximo.

A figura adicional mostra uma configuração distinta dos sites, ilustrando a atualização do diagrama.

\begin{figure}[H]
    \centering
    \begin{subfigure}{.40\linewidth}
        \centering
        \includegraphics[width=\linewidth]{Screenshot 2026-10-24 at 21.48.40.png}
        \caption{Métrica Manhattan.}
        \label{fig:man}
    \end{subfigure}
    \hfill
    \begin{subfigure}{.40\linewidth}
        \centering
        \includegraphics[width=\linewidth]{Screenshot 2026-10-24 at 21.48.47.png}
        \caption{Métrica Chebyshev.}
        \label{fig:cheb}
    \end{subfigure}
    \caption{Comparação das formas das células com três sites e bordas visíveis.}
\end{figure}

\begin{figure}[H]
    \centering
    \includegraphics[width=.53\linewidth]{Screenshot 2026-10-24 at 21.49.01.png}
    \caption{Outra disposição dos três sites, demonstrando a dependência do diagrama em relação às posições dos marcadores.}
\end{figure}

\section{Interação com o utilizador}

A interface permite adicionar um site com um clique numa área vazia, arrastar um marcador existente para alterar a sua posição, deslocar a câmara ao arrastar uma zona livre e ampliar ou reduzir a vista com a roda do rato.

As teclas \texttt{1}, \texttt{2} e \texttt{3} mudam a métrica, \texttt{b} controla as bordas e \texttt{x} remove o site mais próximo da posição do ponteiro.

Como componente extra, foi desenvolvida uma interface em HTML e CSS sobreposta ao \emph{canvas}. O HTML organiza o painel de ajuda, os comandos e o contador de sites, enquanto o CSS define o seu aspeto visual, incluindo transparência, desfoque, tipografia e animação. O painel utiliza \texttt{pointer-events: none}, pelo que permanece visível sem bloquear a interação do utilizador com a aplicação WebGL.

A conversão entre coordenadas do ecrã e do mundo mantém a interação coerente com o pan, o zoom e a proporção da janela.

O diagrama é recalculado no shader sempre que as posições ou os parâmetros de visualização mudam.

\section{Conclusão}

Foi construída uma visualização interativa de diagramas de Voronoi com três métricas e destaque opcional das fronteiras.

As capturas demonstram a associação de cores aos sites e as mudanças geométricas provocadas pela seleção da distância.

Como trabalho futuro, destaca-se a adaptação da espessura das bordas ao zoom, bem como a realização de testes sistemáticos com o limite máximo de sites e diferentes resoluções.

\end{document}
